%!TEX program = xelatex
\documentclass[10pt,a4paper]{article}

\usepackage[a4paper,top=1.18cm,bottom=1.05cm,left=1.45cm,right=1.45cm]{geometry}
\usepackage{xcolor}
\usepackage{fontspec}
\usepackage{tabularx}
\usepackage{enumitem}
\usepackage{titlesec}
\usepackage{hyperref}
\usepackage{microtype}
\usepackage{ragged2e}
\usepackage{array}

\definecolor{msblue}{HTML}{0067B8}
\definecolor{ink}{HTML}{17212B}
\definecolor{muted}{HTML}{4E5968}
\definecolor{rulegray}{HTML}{D6DCE3}

\setmainfont{Avenir Next}
\setsansfont{Avenir Next}
\renewcommand{\familydefault}{\sfdefault}
\linespread{1.04}
\color{ink}

\hypersetup{
  colorlinks=true,
  urlcolor=msblue,
  linkcolor=msblue,
  pdfauthor={Yuwei Zhao},
  pdftitle={Yuwei Zhao - English Resume}
}

\pagestyle{empty}
\setlength{\parindent}{0pt}
\setlength{\parskip}{0pt}
\setlist[itemize]{leftmargin=1.35em,itemsep=2.5pt,topsep=1pt,parsep=0pt,partopsep=0pt}

\titleformat{\section}
  {\large\bfseries\color{msblue}}
  {}{0pt}{}
  [\vspace{-3pt}\color{rulegray}\titlerule]
\titlespacing*{\section}{0pt}{8pt}{10pt}

\newcommand{\datedentry}[3]{%
  \begin{tabularx}{\textwidth}{@{}>{\bfseries}X>{\raggedleft\arraybackslash}p{3.2cm}@{}}
    #1 & \textcolor{muted}{#2}\\[-1pt]
    \multicolumn{2}{@{}>{\raggedright\arraybackslash}p{\textwidth}@{}}{\textcolor{muted}{#3}}
  \end{tabularx}
}

\newcommand{\researchpaperentry}[6]{%
  \item \textbf{#1}\\[-1pt]
  {\small #2\hfill\textbf{\textcolor{msblue}{#3}}}
  \begin{itemize}[leftmargin=1.3em,label={\textcolor{msblue}{\scriptsize\textbullet}},itemsep=1pt,topsep=2pt,parsep=0pt,partopsep=0pt]
    \small
    \item \textbf{Approach:} #4
    \item \textbf{Historical optimization reproduction:} #5
    \item \textbf{Real-world optimization:} #6
  \end{itemize}
}

\newcommand{\paperentrywithsummary}[4]{%
  \item \textbf{#1}\\[-1pt]
  {\small #2\hfill\textbf{\textcolor{msblue}{#3}}}
  \begin{itemize}[leftmargin=1.3em,label={\textcolor{msblue}{\scriptsize\textbullet}},itemsep=1pt,topsep=2pt,parsep=0pt,partopsep=0pt]
    \small
    \item \textbf{Summary:} #4
  \end{itemize}
}

\newcommand{\awardentry}[2]{%
  \noindent #1\hfill\textcolor{muted}{#2}\par\vspace{3.5pt}
}

\begin{document}

\begin{tabularx}{\textwidth}{@{}X>{\raggedleft\arraybackslash}p{10.2cm}@{}}
  {\fontsize{25}{28}\selectfont\bfseries Yuwei Zhao}
  & {\small
      \begin{tabular}{@{}l@{\hspace{1.1em}}l@{}}
        \textbf{Phone:} 173 9804 7360 &
        \textbf{Email:} \href{mailto:zhaoyuwei@stu.pku.edu.cn}{zhaoyuwei@stu.pku.edu.cn}\\[-1pt]
        \multicolumn{2}{@{}l@{}}{\textbf{Website:} \href{https://fyrsta7.github.io/}{https://fyrsta7.github.io/}}
      \end{tabular}
    }\\
\end{tabularx}

\vspace{4pt}
{\color{msblue}\rule{\textwidth}{1.35pt}}\par

\section{Education}
\datedentry{Peking University | School of Computer Science | Ph.D. Student}{2024.09 -- Present}{Advisor: Prof. Yingfei Xiong}
\vspace{4pt}
\datedentry{Peking University | School of Electronics Engineering and Computer Science | B.S.}{2020.09 -- 2024.06}{Major: Information and Computing Science | Thesis: Automatic Synthesis of Dynamic Programming Algorithms via Abstract Interpretation}

\section{Research Profile}
My research focuses on \textbf{LLM-based automated code performance optimization}: using language models to understand programs, identify optimization opportunities, and generate verifiable high-performance implementations. I combine program analysis, synthesis, rule/example guidance, and performance testing, emphasizing gains, reliability, and transferability across languages and large-scale projects.

\section{Publications}
\begin{itemize}
  \researchpaperentry
    {Multi-Source and Cross-Scenario Strategy-Guided Code Optimization}
    {\textbf{Yuwei Zhao}, Qianyu Xiao, Ye Cui, Yijun Yu, Yingfei Xiong}
    {Under Review, 2026}
    {Proposed MoST, a multi-source and cross-scenario strategy-transfer framework that clusters and transfers evidence to generate static-analysis rules, overcoming single-source knowledge and cross-language reuse limitations.}
    {Reproduced developer patches on 151 C/C++, 150 Python, and 50 Rust historical optimization tasks, yielding 21.88\%--180.00\% more successful reproductions than SemOpt.}
    {Discovered and applied previously unseen optimizations in 15 projects, achieving 19.72\%--717.42\% maximum and 4.44\%--258.17\% average improvements. Outperformed SemOpt and Codex on maximum gains in 15/15 projects and average gains in 14/15.}
  \researchpaperentry
    {SemOpt: LLM-Driven Code Optimization via Rule-Based Analysis}
    {\textbf{Yuwei Zhao}, Yuan-An Xiao, Qianyu Xiao, Zhao Zhang, Yingfei Xiong}
    {TOSEM (CCF-A), 2026}
    {Proposed an LLM- and static-analysis-rule-driven optimization framework that distills strategies from historical commits and generates Semgrep rules, addressing the inability of syntax-based retrieval to match semantically equivalent optimizations.}
    {Reproduced developer patches on 151 C/C++ and 150 Python historical optimization tasks. Successful reproductions were 1.38--28$\times$ and 4.60--6.33$\times$ those of the baselines, respectively.}
    {Discovered and applied previously unseen optimizations in 10 large-scale projects, achieving 5.04\%--479.90\% maximum and 1.68\%--143.25\% average improvements.}
  \paperentrywithsummary
    {Superfusion: Eliminating Intermediate Data Structures via Inductive Synthesis}
    {Ruyi Ji, \textbf{Yuwei Zhao}, Nadia Polikarpova, Yingfei Xiong, Zhenjiang Hu}
    {PLDI (CCF-A), 2024}
    {SuFu uses inductive program synthesis to eliminate intermediate data structures in functional programs, solving 264 of 290 tasks.}
  \paperentrywithsummary
    {Decomposition-Based Synthesis for Applying Divide-and-Conquer-Like Algorithmic Paradigms}
    {Ruyi Ji, \textbf{Yuwei Zhao}, Yingfei Xiong, Di Wang, Lu Zhang, Zhenjiang Hu}
    {TOPLAS (CCF-A), 2024}
    {AutoLifter decomposes synthesis through component and variable elimination, automatically applies six divide-and-conquer paradigms, and solves 82 of 96 tasks.}
\end{itemize}

\section{Awards}
\awardentry{Kunpeng Ascend KADC Rising Star Pioneer Award}{2026}
\awardentry{Jiukun Scholarship, School of Computer Science, Peking University}{2025}
\awardentry{Merit Student, Peking University}{2025, 2023, 2021}
\awardentry{Beijing Outstanding Graduate}{2024}
\awardentry{Outstanding Graduate, Peking University}{2024}
\awardentry{Third-Class Scholarship, Peking University}{2023, 2021}
\awardentry{Academic Excellence Award, Peking University}{2022}

\end{document}
